\documentclass[11pt]{article}

\usepackage[T1]{fontenc}
\usepackage[utf8]{inputenc}
\usepackage{amsmath,amssymb,amsthm}
\usepackage{doi}
\usepackage{hyperref}
\usepackage{amsfonts}
\usepackage{amstext}

\newtheorem{theorem}{Theorem}
\newtheorem{lemma}{Lemma}
\newtheorem{corollary}{Corollary}
\newtheorem{proposition}{Proposition}
\theoremstyle{definition}
\newtheorem{definition}{Definition}
\theoremstyle{remark}
\newtheorem{remark}{Remark}

\newcommand{\Heis}{\mathrm{Heis}_3(\mathbb{Z}/q\mathbb{Z})}
\newcommand{\Gal}{\mathrm{Gal}}

\title{Arithmetic Orthogonality of the ADE Case-Selection Gate}
\author{J\'{e}r\^{o}me Beau\\[0.3em]
\small Independent Researcher\\
\small \texttt{jerome.beau@cosmochrony.org}}
\date{2026}

\begin{document}

  \maketitle

  \begin{abstract}
    The generation-split value $\epsilon=\tfrac{1}{10}$ is often presented as a derived prediction of the spectral
    stratigraphy of the corpus.
    This note isolates its exact logical status.
    First, $\epsilon=\tfrac{1}{10}$ is a theorem modulo the ADE case-selection gate: once a case is selected whose
    class-complete spectrum carries the external level ratio $3:2$, the generation-deficit normal form forces
    $\epsilon=\tfrac{1}{10}$ and $\kappa=\tfrac{5}{12}$ algebraically, with no further input.
    Second, the first-principles status of $\epsilon$ is therefore exactly the first-principles status of the gate.
    We then characterise the obstruction to deriving the gate as an arithmetic orthogonality.
    The selection of the $\sqrt{5}$-locus that carries the $3:2$ ratio requires the nontrivial element of
    $\Gal(\mathbb{Q}(\sqrt{5})/\mathbb{Q})$, the action $\sqrt{5}\mapsto-\sqrt{5}$.
    The only field automorphism supplied by the foundations is the Born--Infeld parity $c\mapsto q-c$, which is the
    restriction of complex conjugation $\zeta_q\mapsto\zeta_q^{-1}$ and fixes $\sqrt{5}$.
    In the compositum $K_q=\mathbb{Q}(\zeta_q,\sqrt{5})$ the two actions lie in distinct direct factors of the Galois
    group, hence are arithmetically orthogonal.
    Consequently the gate cannot be discharged from the present axioms without an additional arithmetic principle at
    the spin stratum, unless a type-rigidity property of the recursive non-injectivity tower is established.
    The latter is the single open point that carries the weight of the negative result; we state it and indicate a
    proof direction, but do not conclude it.
  \end{abstract}

  \noindent\textbf{Keywords:} non-injective projection; ADE case selection; binary icosahedral group; golden ratio;
  cyclotomic field; Galois action; Born--Infeld parity; generation split; epistemic stratification.

  \newpage
  \tableofcontents

  \newpage
  \section{Introduction}

    The generation sector of the corpus produces a three-level deficit structure whose outer ratio matches the
    external ADE ratio $3:2$ of a selected case, and this matching fixes the split value $\epsilon=\tfrac{1}{10}$ and
    the common scale $\kappa=\tfrac{5}{12}$~\cite{Beau2026q14,Beau2026a4,Beau2026bim}.
    The value $\epsilon=\tfrac{1}{10}$ has occasionally been described as derived.
    This note draws the line precisely, because the distinction between a posited value and a value forced once a case
    is selected is the whole content of the matter.

    The clarification is in two parts.
    The first part is a trivial decomposition that nonetheless settles the status question: the algebra that takes the
    $3:2$ ratio to $\epsilon=\tfrac{1}{10}$ is general and case-independent, so all of the first-principles content lives
    in the upstream selection of the case~(Section~\ref{sec:conditional}).
    The second part characterises why that upstream selection is not supplied by the present foundations.
    The selection of the $\sqrt{5}$-locus that carries the $3:2$ ratio is an arithmetic act: it requires the Galois
    action that separates the two two-dimensional irreducible characters of the binary icosahedral group $2I$.
    The foundations supply only the Born--Infeld parity, which is complex conjugation on the central character and
    fixes the real quadratic field $\mathbb{Q}(\sqrt{5})$.
    These two actions sit in distinct direct factors of the Galois group of the natural compositum, so the parity
    non-injectivity is not a partial version of the required Galois action; it is orthogonal to
    it~(Sections~\ref{sec:organizing}--\ref{sec:negative}).

    The recursive non-injectivity theorem of ENI~\cite{Beau2026l} guarantees that each stratum of a tower of genuinely
    distinct emergent levels carries its own non-injectivity, but it constrains existence, not arithmetic type.
    The base fibre is closed: its forced involution is the Born--Infeld parity~\cite{Beau2026m}.
    Whether every recursive stratum above the base inherits the parity type, rather than introducing an action in the
    $\sqrt{5}$ factor, is open.
    This type-rigidity statement is the single point that carries the weight of the negative
    result~(Section~\ref{sec:rigidity}).
    Section~\ref{sec:status} records the epistemic stratification and the downstream consequence: observable masses
    route through the explicit operator $E_\Pi$ and the Yukawa sector regardless of how the gate is
    resolved~\cite{Beau2026fmnote,Beau2026prs}.

  \section{The conditional split}
    \label{sec:conditional}

    We work with the generation-deficit normal form of the dynamical spectrum on the generation
    factor~\cite{Beau2026q14}.
    Let $s_0$ be the central eigenvalue, fixed by the projective involution $J_\Pi$, and let $s_\pm$ be the two outer
    eigenvalues, in the normalisation $s_0=1$.

    \begin{definition}[Exit deficit]
For an eigenvalue $s$ with central value $s_0$, the exit deficit is $d_{\mathrm{exit}}(s):=s_0-s$.
The normal form $s_0=1,\ s_\pm=\tfrac{1}{2}\pm\epsilon$ gives
$d_{\mathrm{exit}}(s_0)=0$ and $d_{\mathrm{exit}}(s_\pm)=\tfrac{1}{2}\mp\epsilon$.
    \end{definition}

    \begin{lemma}[Algebraic ratio lemma]
      \label{lem:algebraic}
      If the outer exit deficits are in the ratio $3:2$, that is
      \[
        \frac{d_{\mathrm{exit}}(s_-)}{d_{\mathrm{exit}}(s_+)}=\frac{\tfrac{1}{2}+\epsilon}{\tfrac{1}{2}-\epsilon}=\frac{3}{2},
      \]
      then $\epsilon=\tfrac{1}{10}$.
      If in addition a single common scale $\kappa$ satisfies $\kappa\, d_{\mathrm{exit}}(s_\pm)=|\lambda-1|$ on the two
      outer levels with $\{|\lambda_1-1|,|\lambda_3-1|\}=\{\tfrac{1}{6},\tfrac{1}{4}\}$, then $\kappa=\tfrac{5}{12}$.
    \end{lemma}

    \begin{proof}
      From $2(\tfrac{1}{2}+\epsilon)=3(\tfrac{1}{2}-\epsilon)$ one obtains $1+2\epsilon=\tfrac{3}{2}-3\epsilon$, hence
      $5\epsilon=\tfrac{1}{2}$ and $\epsilon=\tfrac{1}{10}$.
      The outer deficits are then $\{\tfrac{2}{5},\tfrac{3}{5}\}$, and
      $\kappa=(\tfrac{1}{4})/(\tfrac{3}{5})=(\tfrac{1}{6})/(\tfrac{2}{5})=\tfrac{5}{12}$, consistent on both levels.
    \end{proof}

    The lemma is trivial and entirely case-independent: it uses only the ratio $3:2$, not any property of the binary
    icosahedral group or of the generating set.
    The status of $\epsilon=\tfrac{1}{10}$ is therefore the status of the hypothesis that supplies the ratio $3:2$.

    \begin{theorem}[Conditional ADE split]
      \label{thm:conditional}
      Assume the ADE case-selection gate selects the central-lift class over the $\sqrt{5}$-locus, that is the parity
      class over the $A_5$ five-cycle locus, whose class-complete Cayley spectrum has external ratio
      \[
        \frac{|\lambda_3-1|}{|\lambda_1-1|}=\frac{3}{2}
      \]
      with $\{|\lambda_1-1|,|\lambda_3-1|\}=\{\tfrac{1}{6},\tfrac{1}{4}\}$ and central level
      $\lambda_2=1$~\cite{Beau2026a6,Beau2026a4}.
      Then the generation-deficit normal form forces, by Lemma~\ref{lem:algebraic},
      \[
        \frac{\tfrac{1}{2}+\epsilon}{\tfrac{1}{2}-\epsilon}=\frac{3}{2},\qquad
        \epsilon=\frac{1}{10},\qquad \kappa=\frac{5}{12}.
      \]
      The split value is not an independent parameter once the ADE case is selected.
    \end{theorem}

    \begin{remark}[Status]
Theorem~\ref{thm:conditional} is conditional, not predictive.
Its single nontrivial hypothesis is the gate.
The first-principles status of $\epsilon=\tfrac{1}{10}$ is exactly the first-principles status of the gate, and
nothing in the theorem bears on why the $\sqrt{5}$-locus is selected.
The wording avoids committing to ``order five'' versus ``order ten'' for the generating class: the variable-valence
exit analysis uses the order-five generating set~\cite{Beau2026a6}, while the bounded-flux admissibility analysis
uses the order-ten classes $10a\cup 10b$~\cite{Beau2026a1}, and the two are identified by the central lift, as made
precise in Corollary~\ref{cor:tiebreak}.
    \end{remark}

  \section{The arithmetic organizing lemma}
    \label{sec:organizing}

    Two arithmetic data enter the construction from independent parts of the apparatus.
    The root of unity $\zeta_q=e^{2\pi i/q}$ is the central character of the fibre group $\Heis$, fixed by the
    Stone--von Neumann identification of the Weil representation $V_\rho$~\cite{Beau2026m,Beau2026HeisenbergStructure}.
    The quadratic surd $\sqrt{5}$ is the character field of the two two-dimensional irreducible representations of the
    binary icosahedral group $2I$, where it appears through the golden ratio identity
    $\chi(10a)+\chi(10b)=\varphi-\bar{\varphi}=1$ on the order-ten classes~\cite{Beau2026a1,Beau2026a30}.
    These origins are independent: the first is a property of the fibre group on which the axioms act, the second is a
    representation-theoretic property of the spin-stratum subgroup.

    \begin{lemma}[Arithmetic orthogonality]
      \label{lem:organizing}
      Let $q$ be prime with $q\neq 5$, and let $K_q=\mathbb{Q}(\zeta_q,\sqrt{5})$.
      Then $\mathbb{Q}(\zeta_q)$ and $\mathbb{Q}(\sqrt{5})$ are linearly disjoint over $\mathbb{Q}$, so
      \[
        \Gal(K_q/\mathbb{Q})\;\cong\;\Gal(\mathbb{Q}(\zeta_q)/\mathbb{Q})\times\Gal(\mathbb{Q}(\sqrt{5})/\mathbb{Q}).
      \]
      The Born--Infeld parity $c\mapsto q-c$ is the restriction of complex conjugation; as an element of
      $\Gal(K_q/\mathbb{Q})$ it is
      \[
        \text{BI parity}=(\zeta_q\mapsto\zeta_q^{-1},\ \sqrt{5}\mapsto\sqrt{5}).
      \]
      The Galois action separating the two two-dimensional irreducible characters of $2I$ is
      \[
        \text{Galois-spin}=(\zeta_q\mapsto\zeta_q,\ \sqrt{5}\mapsto-\sqrt{5}).
      \]
      These lie in distinct direct factors.
      In particular the Born--Infeld parity fixes $\mathbb{Q}(\sqrt{5})$ pointwise and is not the Galois-spin action.
    \end{lemma}

    \begin{proof}
      The unique quadratic subfield of $\mathbb{Q}(\zeta_q)$ is $\mathbb{Q}(\sqrt{q^{*}})$ with
      $q^{*}=(-1)^{(q-1)/2}q$.
      For prime $q\neq 5$ one has $q^{*}\neq 5$ up to squares, so $\sqrt{5}\notin\mathbb{Q}(\zeta_q)$ and
      $\mathbb{Q}(\zeta_q)\cap\mathbb{Q}(\sqrt{5})=\mathbb{Q}$.
      Linear disjointness gives the direct-product decomposition of the Galois group.
      The central character satisfies
      $e^{2\pi i c/q}\mapsto e^{2\pi i(q-c)/q}=e^{-2\pi i c/q}=\overline{e^{2\pi i c/q}}$ under $c\mapsto q-c$, so the
      Born--Infeld parity is complex conjugation on the $\zeta_q$ factor; since $\sqrt{5}$ is real, complex conjugation
      fixes it.
      The Galois-spin action is the nontrivial element of $\Gal(\mathbb{Q}(\sqrt{5})/\mathbb{Q})$ and fixes $\zeta_q$.
      The two thus occupy distinct direct factors.
    \end{proof}

    \begin{remark}[Why the disjunction is not accidental]
The direct-factor separation reflects the independence of the two origins.
The axioms act on the fibre $\Heis$ and on its cyclotomic factor $\mathbb{Q}(\zeta_q)$; they have no action on the
factor in which $\sqrt{5}$ lives.
The orthogonality is therefore structural, not numerical.
    \end{remark}

  \section{The central-lift tie-break}
    \label{sec:tiebreak}

    The order-ten elements of $2I$ are the order-five elements multiplied by the central element $-I$, since for an
    order-five $g$ one has $(-Ig)^{5}=-I$ and $(-Ig)^{10}=I$.
    The Weil lift of $-I$ is the central reflection $R_b=W(-I)=\mathrm{FT}^{2}$, scalar on the chiral carrier and
    commuting with $\gamma_5$~\cite{Beau2026q11of}.

    \begin{corollary}[Order-five versus order-ten is below parity resolution]
      \label{cor:tiebreak}
      Under the identification $R_b=W(-I)$ of the central lift with the Born--Infeld parity, the projective quotient that
      the construction realises identifies the order-five and order-ten classes rather than distinguishing them.
      The order-five versus order-ten alternative is not a selectable degree of freedom of the framework; it is a
      central-lift ambiguity that the parity erases.
    \end{corollary}

    \begin{proof}
      The pair of classes is exchanged by $\times(-I)$, whose realisation is $R_b$.
      By Lemma~\ref{lem:organizing} the parity acts only on the $\zeta_q$ factor and fixes $\sqrt{5}$, so it permutes the
      classes by the central lift while leaving the Galois data invariant.
      Quotienting by the central lift identifies the two classes; it does not order them.
      Hence the framework supplies no selector between order five and order ten, for the same arithmetic reason that it
      supplies no Galois-spin action: it provides $\times(-I)$, not $\sqrt{5}\mapsto-\sqrt{5}$.
    \end{proof}

    The tie-break and the negative result of Section~\ref{sec:negative} are thus a single fact.
    The parity fixes $\mathbb{Q}(\sqrt{5})$ and permutes only by $\times(-I)$.
    It therefore neither separates order five from order ten, nor separates the two Galois-conjugate characters
    $5a$ and $5b$.

  \section{The negative theorem}
    \label{sec:negative}

    \begin{theorem}[No Galois-spin selection from the present axioms, modulo type-rigidity]
      \label{thm:negative}
      Assume the type-rigidity Lemma~\ref{lem:rigidity}.
      Then every field automorphism forced by the foundations A1--A4 together with the recursive non-injectivity of
      ENI~Corollary~6~\cite{Beau2026l} lies in the factor $\Gal(\mathbb{Q}(\zeta_q)/\mathbb{Q})$ of $\Gal(K_q/\mathbb{Q})$
      and fixes $\mathbb{Q}(\sqrt{5})$ pointwise.
      Consequently the Galois-spin selection, the nontrivial element of $\Gal(\mathbb{Q}(\sqrt{5})/\mathbb{Q})$, is not
      supplied by A1--A4 or by ENI~Corollary~6.
      The $\sqrt{5}$-locus required by the ADE $3:2$ dictionary therefore cannot be selected without an additional
      arithmetic principle at the spin stratum.
    \end{theorem}

    \begin{proof}
      By Lemma~\ref{lem:rigidity} every forced automorphism is a realisation of the Born--Infeld parity, which by
      Lemma~\ref{lem:organizing} lies in $\Gal(\mathbb{Q}(\zeta_q)/\mathbb{Q})$ and fixes $\mathbb{Q}(\sqrt{5})$.
      The Galois-spin action is the nontrivial element of the complementary factor
      $\Gal(\mathbb{Q}(\sqrt{5})/\mathbb{Q})$, which is not in the image of the forced automorphisms.
      Selecting the $\sqrt{5}$-locus over the neutral order-four locus distinguishes the two two-dimensional characters by
      their $\sqrt{5}$ content, hence requires that complementary factor, which the axioms do not provide.
    \end{proof}

    \begin{remark}[Selection against, not merely absence of selection]
The neutral generators derived from Born--Infeld parity and admissibility are those with vanishing spin-$\tfrac{1}{2}$
character, the order-four half-turns; the exact penalised spectral-capacity computation confirms that the order-five
representative of the $\sqrt{5}$-locus is dominated and not selected for $\alpha>0$~\cite{Beau2026a1,Beau2026a2}.
The negative result is therefore an orthogonality between the supplied parity and the required Galois action, not a
mere gap in the derivation.
    \end{remark}

  \section{Type-rigidity along the tower}
    \label{sec:rigidity}

    \begin{lemma}[Type-rigidity, open]
      \label{lem:rigidity}
      Every non-injective projection forced by A1--A4 and ENI~Corollary~6, including the recursive non-injectivities at
      strata above the base fibre, acts on the cyclotomic factor $\mathbb{Q}(\zeta_q)$ and fixes the real quadratic factor
      $\mathbb{Q}(\sqrt{5})$.
    \end{lemma}

    The two regimes of this statement have opposite status, and conflating them is exactly the error to avoid.

    For the base fibre $\Pi_1$ the statement is established.
    The forced fibre involution is the Born--Infeld parity, the minimal fibre condition
    $\Pi^{-1}(y)\supseteq\{\chi,-\chi\}$ closed by A3 and the Stone--von Neumann identification of
    $V_\rho$~\cite{Beau2026m,Beau2026HeisenbergStructure}.
    By Lemma~\ref{lem:organizing} this involution lies in $\Gal(\mathbb{Q}(\zeta_q)/\mathbb{Q})$ and fixes $\sqrt{5}$.

    For the recursive strata $\Pi_{i\geq 2}$ the statement is open.
    ENI~Corollary~6 forces $S_{\Pi_i}>0$ at each stratum of a tower of genuinely distinct emergent levels, but it
    constrains existence of the non-injectivity, not its arithmetic type~\cite{Beau2026l,Beau2026ninote}.
    Nothing in the corollary forbids a stratum from realising its non-injectivity by an automorphism in the $\sqrt{5}$
    factor.
    The spin stratum is precisely such a recursive stratum, and it is where $\sqrt{5}$ enters through the representation
    theory of $2I$.

    \begin{remark}[Proof direction, not concluded]
The entropy monotonicity $S_{\Pi_2\circ\Pi_1}\geq S_{\Pi_1}$ of ENI~Corollary~6 constrains the composition of
strata.
If the only automorphism available to realise a fibre non-injectivity is the parity, a candidate argument is that
the composition cannot create an arithmetic factor absent from the individual stages, so that the recursive type is
inherited from the base.
This is a type-inheritance claim, distinct from entropy preservation, and the step from one to the other is not
filled.
Lemma~\ref{lem:rigidity} is therefore stated as structural and open, and Theorem~\ref{thm:negative} is conditional on
it.
    \end{remark}

  \section{Status and downstream consequence}
    \label{sec:status}

    The epistemic stratification of the six statements is as follows.

    \begin{center}
      \begin{tabular}{lll}
        \hline
        Statement & Content & Status\\
        \hline
        Lemma~\ref{lem:algebraic} & ratio $3:2\Rightarrow\epsilon=\tfrac{1}{10}$ & proved, trivial, case-independent\\
        Theorem~\ref{thm:conditional} & gate $\Rightarrow$ ratio $3:2\Rightarrow\epsilon$ & conditional / open hypothesis\\
        Lemma~\ref{lem:organizing} & parity $\perp$ Galois-spin in $\Gal(K_q/\mathbb{Q})$ & proved\\
        Corollary~\ref{cor:tiebreak} & order five $=$ order ten mod central lift & proved, modulo $R_b=W(-I)$\\
        Theorem~\ref{thm:negative} & no Galois-spin from A1--A4 $+$ ENI Cor.~6 & proved modulo Lemma~\ref{lem:rigidity}\\
        Lemma~\ref{lem:rigidity} & type-rigidity along the tower & structural / open\\
        \hline
      \end{tabular}
    \end{center}

    The net statement is that $\epsilon=\tfrac{1}{10}$ is a theorem modulo the ADE case-selection gate, that the gate's
    first-principles status reduces to the type-rigidity Lemma~\ref{lem:rigidity}, and that this reduction is an
    arithmetic orthogonality rather than an unresolved computation.
    This converts the obstruction from ``unproven'' to ``characterised as requiring an arithmetic principle foreign to
    A1--A4 and ENI~Corollary~6, unless type-rigidity holds''.

    The downstream consequence bounds the relevance of the gate.
    The gate fixes a ratio of levels and hence the structural input $\epsilon$; it does not fix observable masses.
    The charged-lepton and quark spectra and the mixing matrices require the explicit construction of the operator
    $E_\Pi$ and the projected Yukawa sector, with the mass hierarchy carried by the cascade amplification and its
    exponent~\cite{Beau2026fmnote,Beau2026prs,Beau2026a4}.
    Even a derivation of the gate would not by itself produce a mass.
    The gate is therefore a finite, characterisable structural sub-problem that can be held at conditional status while
    the quantitative effort proceeds on $E_\Pi$, the explicit Lorentzian block, and the Yukawa sector.

    \newpage
    \bibliographystyle{plain}
    \bibliography{cosmochrony-bibliography}

\end{document}
